\chapter*{Resumen}

Al día de hoy, las redes sociales han transformado la manera en que las personas se comunican e interactúan entre sí, además de ser un espacio libre para compartir información, noticias, fotos y videos; sin embargo, la libre expresión y falta de censura, esto sumado a la sensación de anonimato por parte de los usuarios, ha resultado en un aumento de conductas negativas y perjudiciales para el internauta, tales como el acoso, la discriminación y la intimidación. Por ello, el presente proyecto propone el diseño y desarrollo de un sistema de detección automática de lenguaje ofensivo que pueda detectar dichos mensajes de odio para proteger al usuario.

El objetivo principal es ofrecer una solución precisa, adaptable y funcional en tiempo real, dirigida inicialmente a medios de comunicación digitales que enfrentan limitaciones para abrir sus espacios de comentarios por temor a contenidos ofensivos. El sistema podrá integrarse como servicio externo o mediante API, permitiendo filtrar contenido antes de su publicación pública, incluso en escenarios de transmisión en vivo. El proyecto considera el diseño del pipeline de clasificación, la preparación y limpieza del dataset, el entrenamiento y evaluación del modelo, así como una proyección de viabilidad económica, escalabilidad técnica y análisis de mercado.

Como resultado, se espera contribuir al desarrollo de herramientas inclusivas y accesibles para la moderación automática del lenguaje ofensivo en español, promoviendo entornos digitales más seguros y abiertos a la participación constructiva del público.

\noindent\textbf{Palabras clave:} Lenguaje ofensivo, procesamiento del lenguaje natural, moderación automática de contenido, transformers
